% p.58 — الباب الرابع والعشرون في ما يلحق مع الفعل واسمي الفاعل والمفعول
% Layout: six vertical lists of 21 cells each, printed as six columns of one grid.
% Sub-headings occur mid-column (col. 2: «مضارع مرفوع» after 13 cells; col. 3: «مضارع مجزوم»
% after 13; col. 4: «مُضَارِعٌ مَنْصُوْبٌ» after 6; col. 5 is headed «اسم فاعل و مفعول»;
% col. 6 has «هُوَ فَاعِلٌ» ruled off as a heading after 3 cells). Column 1 is headed «ماضي»;
% column 2 has no heading (its first cell إِذْ فَعَلَ sits in the heading row).
% The partial rules around those sub-headings are not reproduced. Columns are read top to
% bottom, right to left. Initial hamza of particles is printed as bare alif + vowel in the
% lithograph (اِنْ، اِذَا); transcribed with the hamza letter.
\mspage{58}
\bab{الباب الرابع والعشرون في ما يلحق مع الفعل واسمي الفاعل والمفعول}
\begin{sarf}{|c|c|c|c|c|c|}
\thead{ماضي} & إِذْ فَعَلَ & هَلَّا يَفْعَلُ & إِنْ لَا يَفْعَلْ & \thead{اسم فاعل و مفعول} & مَعَ أَنَّهُ فَاعِلٌ\\\hline
هَلْ فَعَلَ & إِذًا فَعَلَ & أَمَا يَفْعَلُ & إِنْ لَمْ يَفْعَلْ & هُوَ فَاعِلٌ & حَتّٰى أَنَّهُ فَاعِلٌ\\
قَدْ فَعَلَ & حَتّٰى فَعَلَ & أَلَا يَفْعَلُ & حَتّٰى إِذَا لَمْ يَفْعَلْ & إِنَّهُ فَاعِلٌ & أَإِنَّهُ فَاعِلٌ\\
مَا فَعَلَ & حَتّٰى إِذَا فَعَلَ & إِذْ يَفْعَلُ & \unclear{كَانَ} لَمْ يَفْعَلْ & لٰكِنَّهُ فَاعِلٌ & \thead{هُوَ فَاعِلٌ}\\
لَا فَعَلَ & لَمَّا فَعَلَ & إِذَا يَفْعَلُ & لَوْ لَمْ يَفْعَلْ & أَنَّهُ فَاعِلٌ & هُمَا فَاعِلَانِ\\
آفَعَلَ & أَنْ فَعَلَ & إِذَا مَا يَفْعَلُ & أَيٌّ يَفْعَلْ & لَعَلَّهُ فَاعِلٌ & هُمْ فَاعِلُوْنَ\\
مَاذَا فَعَلَ & لِأَنْ فَعَلَ & قَدْ يَفْعَلُ & \thead{مُضَارِعٌ مَنْصُوْبٌ} & لَيْتَهُ فَاعِلٌ & هِيَ فَاعِلَةٌ\\
مَتٰى فَعَلَ & عَلٰى أَنْ فَعَلَ & قَدْ لَا يَفْعَلُ & لَنْ يَفْعَلَ & إِنَّمَا هُوَ فَاعِلٌ & هُمَا فَاعِلَتَانِ\\
أَيْنَ فَعَلَ & مَا فَعَلَ & وَلٰكِنْ يَفْعَلُ & كَيْ يَفْعَلَ & أَنَّمَا هُوَ فَاعِلٌ & هُنَّ فَاعِلَاتٌ\\
أَنّٰى فَعَلَ & عَلٰى مَا فَعَلَ & كَمَا يَفْعَلُ & إِذَنْ يَفْعَلَ & لٰكِنَّمَا هُوَ فَاعِلٌ & أَنْتَ فَاعِلٌ\\
لِمَ فَعَلَ & رُبَّمَا فَعَلَ & رُبَّمَا يَفْعَلُ & أَنْ يَفْعَلَ & لَيْتَمَا هُوَ فَاعِلٌ & أَنْتُمَا فَاعِلَانِ\\
لِمَاذَا فَعَلَ & كَمَا فَعَلَ & لَوْلَا يَفْعَلُ & لِأَنْ يَفْعَلَ & لَعَلَّمَا هُوَ فَاعِلٌ & أَنْتُمْ فَاعِلُوْنَ\\
أَمَا فَعَلَ & لَمَا فَعَلَ & لَيَفْعَلُ & حَتّٰى يَفْعَلَ & كَأَنَّمَا هُوَ فَاعِلٌ & أَنْتِ فَاعِلَةٌ\\
أَلَا فَعَلَ & \thead{مضارع مرفوع} & \thead{مضارع مجزوم} & عَلٰى أَنْ يَفْعَلَ & إِذَا مَا هُوَ فَاعِلٌ & أَنْتُمَا فَاعِلَتَانِ\\
كُلَّمَا فَعَلَ & سَيَفْعَلُ & لَا يَفْعَلْ & عَلٰى أَنْ لَا يَفْعَلَ & لِمَا هُوَ فَاعِلٌ & أَنْتُنَّ فَاعِلَاتٌ\\
لَعَلَّمَا فَعَلَ & سَوْفَ يَفْعَلُ & لِيَفْعَلْ & لِأَنْ لَا يَفْعَلَ & كَمَا هُوَ فَاعِلٌ & أَنَا فَاعِلٌ\\
لَيْتَمَا فَعَلَ & لَا يَفْعَلُ & لَمْ يَفْعَلْ & كَيْ لَا يَفْعَلَ & هَلْ هُوَ فَاعِلٌ & أَنَا فَاعِلَةٌ\\
إِذَا مَا فَعَلَ & هَلْ يَفْعَلُ & لَمَّا يَفْعَلْ & إِلٰى أَنْ يَفْعَلَ & أَهُوَ فَاعِلٌ & نَحْنُ فَاعِلَانِ\\
إِنْ فَعَلَ & أَيَفْعَلُ & إِنْ يَفْعَلْ & إِلٰى أَنْ لَا يَفْعَلَ & رُبَّهُ فَاعِلٌ & نَحْنُ فَاعِلَتَانِ\\
إِذَا فَعَلَ & مَاذَا يَفْعَلُ & مَنْ يَفْعَلْ & لِكَيْ لَا يَفْعَلَ & لِأَنَّهُ فَاعِلٌ & نَحْنُ فَاعِلُوْنَ\\
لَوْ فَعَلَ & لِمَ يَفْعَلُ & مَا يَفْعَلْ & لِيَفْعَلَ & عَلٰى أَنَّهُ فَاعِلٌ & نَحْنُ فَاعِلَاتٌ\\
\end{sarf}
